\writeSectionFrame{\presentationTitle}{Fazit}
\begin{frame}{Anwendbarkeit}
	\begin{itemize}
 		\item Remote-Proxy
 		\item Protection-Proxy
 		\item Cache-Proxy
 		\item Synchronization-Proxy
 		\item Virtual-Proxy
 		\item Counting-Proxy
 		\item Firewall-Proxy
 	\end{itemize}
\end{frame}

\note{
	Das Proxy-Muster kann sehr vielfältig eingesetzt werden. Ein Remote-Proxy ist ein lokaler Vertreter für ein
	Objekt auf einem anderen Rechner. Ein Protection-Proxy schützt Komponenten vor unerlaubten Zugriffen.
	Über einen Cache-Proxy können sich lokale Klienten die Ergebnisse entfernter Komponenten teilen. Ein 
	Synchronization-Proxy erlaubt die Synchronisation mehrerer gleichzeitiger Zugriffe auf eine Komponente.
	Ein Virtual-Proxy erzeugt teure Objekte auf Verlangen. Ein Counting-Proxy dient als Ersatz für einen einfachen
	Zeiger, der zusätzlich folgende Funktionen anbietet: Zählen der Referenzen auf das eigentliche Objekt und  
	automatische Freigabe, wenn es keine Referenzen mehr besitzt. Ein Firewall-Proxy schützt lokale Klienten vor Zugriffen
	von außen. Dabei muss das Proxy-Muster auch nicht immer im Zusammenhang mit objektorientierten Programmen verstanden
	werden, sondern kann sehr viel vielfältiger zum Einsatz kommen. 
}

\begin{frame}{Weitere Aufgaben}
	\begin{itemize}
 		\item Serialisierung
 		\item Austauschen von Objekten zur Laufzeit
 		\item Abbildung von Aspekten
 		\item Konvertierung von Datenformaten
 		\item Änderung von Funktionalitäten
        \item Ressourcenmanagment
        \item Lazy Loading
        \item Zugriffskontrolle
 	\end{itemize}
\end{frame}

\note{
	Proxies können darüber hinaus noch viele weitere sinnvolle Aufgaben wahrnehmen. Sie können Daten serialisieren.	Objekte können zur Laufzeit ausgetauscht werden, was z.B. zu Wartungsgründen passieren kann. Sie können auch sog. Aspekte abbilden, also Funktionalitäten, die nichts mit der eigentlichen Geschäftslogik zu tun haben wie	z.B. Logging, die nicht in das echte Subjekt eingebaut werden sollen. Auch eine Konvertierung von Datenformaten wie bspw. von ANSI nach Unicode bei Strings wäre denkbar. 
}

\realworld{Internet-Proxies}{Vermittler zwischen einem Benutzer und dem Internet}
\note{
    Ein Internet-Proxy-Server fungiert als Vermittler zwischen einem Benutzer und dem Internet. Er kann Anfragen von Benutzern entgegennehmen und sie im Namen der Benutzer an das Internet weiterleiten. Der Proxy kann die Identität der Benutzer verbergen. Häufig besuchte Webseiten können zwischengespeichert werden, um die Ladezeiten zu verkürzen und die Bandbreite zu sparen. Organisationen können den Zugang zu bestimmten Webseiten oder Online-Diensten einschränken.
}

\realworld{Datenbank-Proxies}{Kommunikation zwischen einer Anwendung und der Datenbank}
\note{
    Ein Datenbank-Proxy verwaltet die Kommunikation zwischen einer Anwendung und der Datenbank. Er kann Funktionen wie Caching, Verbindungspooling und Lastverteilung übernehmen.
}

\realworld{Content Delivery Networks (CDNs)}{Geografisch verteilte Webserver}
\note{
    Ein CDN verwendet Proxy-Server, die geografisch verteilt sind, um Webinhalte schneller an Benutzer zu liefern. Die Inhalte werden an Server in der Nähe des Benutzers zwischengespeichert. Inhalte werden von einem Server bereitgestellt, der näher am Benutzer liegt. Der Traffic wird auf viele Server verteilt, um Überlastungen zu vermeiden. Bei Ausfall eines Servers kann ein anderer Server die Anfragen übernehmen.
}

\realworld{Remote-Proxies}{Zugriff auf einen entfernten Dienst}
\note{
    Remote-Proxies werden oft in verteilten Systemen verwendet, um einen Zugriff auf ein entferntes Objekt oder einen Dienst zu ermöglichen, als wäre es lokal. Der Proxy kann den Zugriff auf das entfernte Objekt steuern und der Proxy kann Netzwerkübertragungen optimieren und Daten komprimieren.
}

\realworld{Server-Side Proxies für APIs}{Weiterleitung von Client-Anfragen an externe APIs}
\note{
    In Webanwendungen können Server-Side Proxies verwendet werden, um Anfragen von Clients an externe APIs weiterzuleiten und zusätzliche Funktionen wie Authentifizierung, Rate Limiting oder Caching zu integrieren. Dadurch wird z.B. verhindert, dass API-Schlüssel direkt im Client-Code eingebettet werden und die Anzahl der Anfragen an die externe API kann gesteuert werden, um Überlastung zu verhindern.
}

\realworld{ORM (Object-Relational Mapping) Systeme}{Lazy Loading}
\note{
    ORMs wie Hibernate in Java oder Entity Framework in .NET nutzen Proxys, um Lazy Loading von Daten zu unterstützen. Wenn ein ORM-Proxy erstellt wird, wird die Datenbankabfrage erst ausgeführt, wenn die Daten tatsächlich benötigt werden. Dies reduziert die Anzahl der Datenbankabfragen und verbessert die Performance durch verzögertes Laden.
}

\vorteil{Kontrolle über einen Dienst, ohne dass die Clients etwas davon wissen}
\note{
    Da der Client nur die Schnittstelle kennt und nicht weiß, ob er gerade mit dem realen Objekt oder dem Proxy kommuniziert, hat man volle Kontrolle über den angebotenen Dienst. Der Client bekommt nichts davon mit, wenn das Proxy-Objekt gegen das reale Objekt ausgetauscht wird.
}

\vorteil{Unabhängigkeit vom eigentlichen Dienst}
\note{
    Wie bei dem WebService-Beispiel angedeutet, kann der Proxy dazu dienen, einen nicht funktionierenden Dienst zu ersetzen.
}

\vorteil{Open-Closed-Prinzip}
 \note{
    Neue Proxies können einfach eingeführt werden. 
 }

\vorteil{Vielfältig einsetzbar}
\note{
    Wie bereits oben gezeigt, ist das Proxymuster sehr vielfältig einsetzbar.
}

\vorteil{Performance-Vorteile}
\note{
    Wenn man einen Proxy für Dinge wie z.B. Caching oder Lazy Loading verwendet, kann er dabei helfen, große Performance-Verbesserungen zu erreichen.
}

\nachteil{Erhöhte Komplexität}
\note{
    Der Einsatz eines Proxys fügt eine zusätzliche Schicht zwischen dem Client und dem realen Objekt hinzu, was die Systemkomplexität erhöht. Die Implementierung und Wartung von Proxys kann zusätzlichen Code und Testing-Aufwand erfordern.
}

\nachteil{Leistungsengpässe möglich}
\note{
    Obwohl der Proxy Vorteile in Bezug auf Lazy Loading und Caching bieten kann, kann er auch zusätzlichen Overhead verursachen, da jede Anfrage durch den Proxy geleitet wird. In Systemen mit sehr hoher Nachfrage kann der Proxy selbst zu einem Engpass werden, wenn er nicht effizient implementiert ist.
}

\nachteil{Potenzielle Designprobleme}
\note{
    Ein unsachgemäß implementierter Proxy kann Designprobleme wie hohe Kopplung oder unnötige Komplexität verursachen. Beispiel: Ein Proxy, der zu viele Verantwortlichkeiten übernimmt oder nicht richtig konfiguriert ist, kann das System schwerer wartbar machen.
}

\nachteil{Potenzielle Redundanz}
\note{
    Wenn nicht sorgfältig implementiert, kann der Proxy redundant sein oder zusätzliche Logik hinzufügen, die nicht unbedingt benötigt wird.
}